\documentclass[12pt]{extarticle}
\usepackage{array}
\usepackage[T1]{fontenc}
\usepackage[Russian]{babel}
\usepackage{booktabs}
\usepackage{graphicx}
\usepackage{spreadtab}
\usepackage{numprint}
\usepackage[margin=3cm]{geometry}
\usepackage{float}
\usepackage{xfp}
\usepackage{amsmath} 
\usepackage{indentfirst}  
\begin{document}
\author{Осипов Александр Б06-604}
\title{Лабораторная работа 1.4.5. Изучения колебаний струны}
\maketitle

\section{Цель}

изучить поперечные стоячие волн на тонкой натянутой струне;
измерить собственные частоты колебаний струны и проверить условие образования
стоячих волн; измерить скорость распространения поперечных волн на струне и исследовать её зависимость от натяжения струны.

\vspace{5mm}

\section{Оборудование}

закрепленная на станине стальная струна, набор грузов, электромагнитные датчики, звуковой генератор, двухканальный осциллограф,
частотомер.

\vspace{5mm}

\section{Ход работы}



\end{document}